\tutorial{6}{Calibration with Ostrich (DDS)}

\begin{goal}
You calibrate Tutorial Creek against streamflow at the outlet and heads at three observation wells, with OSTRICH's
Dynamically Dimensioned Search (DDS). The observations were made with known (``true'') parameters, so you can check what
the calibration recovers. First you ask which parameters the observations can see at all, and look at a first attempt
that ran into parameters they cannot see or cannot tell apart; then you calibrate the two that they see clearly and
separately --- the hydraulic conductivities of the valley sand and of the slope till --- and validate the result on two
years it has not seen. The setup follows the official Raven--Ostrich tutorial (Demo~C1:
\cmd{ostIn.txt}, a template, \cmd{Ost-RAVEN.sh} and \cmd{save\_best.sh}). Time: about \V{t6:ddsmin}~minutes for the
calibration (\V{t6:maxiter} model runs) and 30 minutes of work.
\end{goal}

\section{Get OSTRICH}
OSTRICH (L.~Shawn Matott) is not part of the package. Get it from its official sources:
\begin{itemize}
\item \url{https://www.civil.uwaterloo.ca/envmodelling/Ostrich.html} (University of Waterloo; Linux, macOS and Windows
      executables, and the Raven--Ostrich tutorial files with Demo~C1);
\item \url{https://www.eng.buffalo.edu/~lsmatott/Ostrich/OstrichMain.html} (the author's page; source code, \cmd{make GCC});
\item \url{https://usbr.github.io/ostrich} (documentation of the current versions).
\end{itemize}
This tutorial was made with \cmd{Ostrich\_v20171219\_GCC\_Linux.exe} of the Raven--Ostrich tutorial files (Demo~C1), which
reports itself as \emph{OSTRICH version \V{t6:ostver}}. Unzip the Raven--Ostrich tutorial files, copy
\cmd{Demo\_C1/Ostrich\_v20171219\_GCC\_Linux.exe} into a folder \cmd{\textasciitilde/ostrich}, and put its full path in a
variable:
\begin{shell}
$ mkdir -p ~/ostrich
$ cp Demo_C1/Ostrich_v20171219_GCC_Linux.exe ~/ostrich/
$ export OSTRICH=~/ostrich/Ostrich_v20171219_GCC_Linux.exe
$ chmod +x $OSTRICH
\end{shell}
\begin{trouble}
\cmd{cannot execute: required file not found} (although the file is there): that executable was built on a cluster and
names a program loader that your Linux does not have. Start it through your system's loader instead:
\cmd{export OSTRICH="/lib64/ld-linux-x86-64.so.2 \$HOME/ostrich/Ostrich\_v20171219\_GCC\_Linux.exe"} (this is how it ran
for this tutorial), or compile OSTRICH from its source.
\end{trouble}

\section{The calibration folder}
\begin{shell}
$ cp -r $TUT/t6_calibration/start ~/work/t6
$ cd ~/work/t6
\end{shell}
As in Demo~C1, the folder holds the OSTRICH files and the Raven model in \cmd{model/}:
\begin{center}\small
\begin{tabularx}{\linewidth}{>{\raggedright\arraybackslash}p{0.27\linewidth}X}\toprule
File & Role\\\midrule
\cmd{ostIn.txt} & OSTRICH's setup: the algorithm, the parameters, the responses and the objective\\
\cmd{Creek.rvp.tpl} & template of \cmd{Creek.rvp}: the parameters as names, which OSTRICH replaces by values\\
\cmd{Ost-RAVEN.sh} & runs the model once: copies the new \cmd{Creek.rvp} into \cmd{model/}, runs Raven\\
\cmd{save\_best.sh} & saves the files of the best run so far into \cmd{best/}\\
\cmd{model/} & the model of Tutorial~2 for 2010--2012, with the initial parameter guess and the observations\\\bottomrule
\end{tabularx}\end{center}

\subsection{The observations}
\cmd{model/CreekObs.rvt} holds daily streamflow at the outlet (subbasin~2) and weekly heads at the three observation wells of
Tutorial~2, for 2010--2014. They were made by running the model of Tutorial~2 --- whose parameters are the ``true'' ones
of this experiment --- and adding noise (\V{t6:noiseq}\% on the flows, \V{t6:noiseh}~cm on the heads):
\term{t6_obs_head.txt}
\cmd{:ObservationData HYDROGRAPH 2 m3/s} is streamflow at subbasin~2; \cmd{:ObservationData GW\_HEAD 101 m} heads at the
observation well with ID~101 of \cmd{:ObservationWells} (Raven compares an observed head on date D with the simulated head at
D); \cmd{-1.2345} marks a day without a value. Each series has an \cmd{:ObservationWeights} block of the same length: 0
in 2010, the model's warm-up year, and 1 afterwards, so the first year does not count in \cmd{Diagnostics.csv}. The model
reads the file through its \cmd{.rvt}:
\diff{t6_rvt.diff}

\subsection{The initial guess}
\cmd{model/Creek.rvp} has the parameters of a first guess, not the true ones. Run it once by hand:
\begin{shell}
$ cd model
\end{shell}
\term{t6_run_init.txt}
\term{t6_diag_init.txt}
One line per observed series: \cmd{HYDROGRAPH\_ALL[2]} and \cmd{GW\_HEAD\_ALL[101]} to \cmd{[103]}, in the order of the
\cmd{.rvt}, with the metrics of \cmd{:EvaluationMetrics} in \cmd{Creek.rvi}. The first guess has the sand and the till
conductivities wrong (the other parameters are those of Tutorial~2). The flow is not bad
(KGE \V{t6:kgeInit}), but the heads are off by \V{t6:rmseInit101}~m at OW\_NEAR and \V{t6:rmseInit103}~m at OW\_SLOPE
(RMSE). Go back up:
\begin{shell}
$ cd ..
\end{shell}

\section{Which parameters can the observations see?}\label{sec:t6ident}
Before calibrating, check that the observations depend on the parameters at all. The table below changes one parameter at a
time by 10\% and shows how the KGE of the streamflow, the head RMSE of each well and the objective of the calibration
change. The objective is defined in \cmd{ostIn.txt} below: $-\mathrm{KGE}$ plus, for each well, its head RMSE times
$1/(2\sigma)$ of its observed heads (\V{t6:w101}, \V{t6:w102} and \V{t6:w103} for OW\_NEAR, OW\_UPPER and OW\_SLOPE).
With the true parameters the KGE is \V{t6:sens:kge}, the RMSEs are \V{t6:sens:rmse101}, \V{t6:sens:rmse102} and
\V{t6:sens:rmse103}~mm (the noise of the observations) and the objective is \V{t6:sens:base}.
\begin{center}\small
\begin{tabular}{llrrrrr}\toprule
 & & & \multicolumn{3}{c}{change of the head RMSE [mm]} & \\\cmidrule(lr){4-6}
Parameter & change & KGE & OW\_NEAR & OW\_UPPER & OW\_SLOPE & objective\\\midrule
$K_h$ GRAVEL & $-$10\% & $-$0.006 & $-$1 & 0 & +1 & +0.002\\
$K_h$ GRAVEL & +10\% & $-$0.005 & +6 & +4 & +1 & +0.072\\[2pt]
$S_y$ GRAVEL & $-$10\% & 0.000 & 0 & 0 & 0 & +0.004\\
$S_y$ GRAVEL & +10\% & 0.000 & 0 & 0 & 0 & $-$0.001\\[2pt]
$K_h$ SAND & $-$10\% & $-$0.002 & $-$1 & 0 & +19 & +0.003\\
$K_h$ SAND & +10\% & $-$0.002 & +1 & +1 & +15 & +0.023\\[2pt]
$K_h$ TILL & $-$10\% & 0.000 & 0 & 0 & +1\,845 & +0.776\\
$K_h$ TILL & +10\% & 0.000 & +1 & 0 & +1\,598 & +0.674\\[2pt]
\cmd{PERC\_COEFF} & $-$10\% & $-$0.016 & +2 & +1 & +220 & +0.122\\
\cmd{PERC\_COEFF} & +10\% & $-$0.017 & 0 & 0 & +185 & +0.090\\
\bottomrule
\end{tabular}

\end{center}
(True values: $K_h$ of GRAVEL \V{t6:trueK_GRAVEL}~m/d, valley, layer~1; $S_y$ of GRAVEL \V{t6:trueSY_GRAVEL}; $K_h$ of SAND
\V{t6:trueK_SAND}~m/d, valley, layer~2; $K_h$ of TILL \V{t6:trueK_TILL}~m/d, slopes; \cmd{PERC\_COEFF} of SUBSOIL
\V{t6:truePERC_SUBSOIL}~1/d.)
\begin{itemize}
\item The gravel is seen by the valley wells: OW\_NEAR and OW\_UPPER respond to its conductivity, more to a higher value
      than to a lower one, and with their large weights that is most of its change of the objective. Its $S_y$ hardly shows.
\item The sand is seen only at OW\_SLOPE, where it carries the groundwater of the slopes down the valley: the RMSE there
      grows by \V{t6:sens:sandslopemin} to \V{t6:sens:sandslopemax}~cm, about the noise, and the small weight of OW\_SLOPE
      (\V{t6:w103}, set by the large spread of its heads) cuts that to hundredths of the objective. The valley wells hardly
      see it, although OW\_NEAR is screened in the sand 200~m from PW\_1: its heads follow the creek and the gravel.
\item The till matters enormously: the water table under the slopes lies at a height set by the ratio of recharge to
      conductivity, and a 10\% change of its $K_h$ raises the RMSE at OW\_SLOPE to \V{t6:sens:K_TILL:lo:slope} and
      \V{t6:sens:K_TILL:hi:slope}~m.
\item \cmd{PERC\_COEFF} sets how much of the subsoil's drainage becomes recharge (the rest flows to the creek as interflow),
      so it shows in the streamflow and, through the recharge, at OW\_SLOPE too.
\end{itemize}
The objective is not exactly at its minimum at the true values: the observations are noisy, so a neighbouring set can fit
them a little better (a negative change in the table). Read the table for the size of the changes on either side.

\subsection{A first attempt}
A first calibration, kept in \cmd{\$TUT/t6\_calibration/finished/first\_attempt/}, searched four parameters --- $K_h$ and
$S_y$ of the gravel, $K_h$ of the till and \cmd{PERC\_COEFF} --- with \V{t6:fa:maxiter} DDS runs:
\term{t6_fa_params.txt}
\begin{center}\small
\begin{tabular}{lrrrr}\toprule
Parameter & initial & true & found & error\\\midrule
$K_h$ of GRAVEL [m/d] & \V{t6:fa:initK_GRAVEL} & \V{t6:trueK_GRAVEL} & \V{t6:fa:bestK_GRAVEL} & \V{t6:fa:errK_GRAVEL}\%\\
$S_y$ of GRAVEL & \V{t6:fa:initSY_GRAVEL} & \V{t6:trueSY_GRAVEL} & \V{t6:fa:bestSY_GRAVEL} & \V{t6:fa:errSY_GRAVEL}\%\\
$K_h$ of TILL [m/d] & \V{t6:fa:initK_TILL} & \V{t6:trueK_TILL} & \V{t6:fa:bestK_TILL} & \V{t6:fa:errK_TILL}\%\\
\cmd{PERC\_COEFF} of SUBSOIL [1/d] & \V{t6:fa:initPERC_SUBSOIL} & \V{t6:truePERC_SUBSOIL} & \V{t6:fa:bestPERC_SUBSOIL} & \V{t6:fa:errPERC_SUBSOIL}\%\\\bottomrule
\end{tabular}\end{center}
Its objective (which weighted every well alike, \V{t6:headw}~times its head RMSE) reached \V{t6:fa:cost} at run
\V{t6:fa:itBest}; with the true parameters it is \V{t6:fa:costtrue}. Two things went wrong, both foreseen by the table above:
\begin{itemize}
\item The gravel's parameters ended where the objective is flat; no number of runs pins them down with these observations.
\item DDS found a till conductivity that fitted OW\_SLOPE early (it kept the value of run \V{t6:fa:ktfrom} to the end),
      but together with a wrong \cmd{PERC\_COEFF}: less recharge and a slightly lower conductivity give the same heads on
      the slope. From there, changing either parameter alone costs much more (the table) than the streamflow gains, so
      the search stays in that narrow valley, even over \V{t6:fa:maxiter} runs.
\end{itemize}
Two changes make the calibration well posed.
\begin{enumerate}
\item Calibrate what the data can see, and not two parameters that they see only together. The calibration below searches
      the hydraulic conductivity of the till, which the data see clearly, and that of the sand, which they see weakly ---
      to show what becomes of such a parameter. The gravel and \cmd{PERC\_COEFF} keep the values of
      Tutorial~2, as if known from elsewhere (a pumping test, a calibration of the streamflow alone); in this synthetic
      experiment those are their true values, so the test is whether DDS recovers the two conductivities.
\item Weight each well by the spread of its own observations. In the calibration years OW\_SLOPE's heads range over
      \V{t6:swingm103}~m, those of OW\_NEAR and OW\_UPPER over \V{t6:swing101} and \V{t6:swing102}~cm. Weighted alike, as in the first attempt, OW\_SLOPE dominates the objective so much
      that DDS keeps any set that fits it, and the sand and the till also trade off there (a higher sand conductivity drains
      the slopes, a lower till conductivity raises them again). Measured against its own spread, each well and the
      streamflow get a comparable say.
\end{enumerate}

\subsection{The template}
\cmd{Creek.rvp.tpl} is \cmd{model/Creek.rvp} with two values replaced by parameter names:
\diff{t6_tpl.diff}

\subsection{\texttt{ostIn.txt}}
\showfile{t6_calibration/start/ostIn.txt}
\begin{itemize}
\item \cmd{ProgramType DDS}, \cmd{ObjectiveFunction GCOP}: DDS minimises a general objective built from the responses.
\item \cmd{ModelExecutable} and \cmd{PreserveBestModel}: the two scripts below.
\item \cmd{BeginFilePairs}: OSTRICH writes \cmd{Creek.rvp} from \cmd{Creek.rvp.tpl} before each run.
\item \cmd{BeginParams}: initial value, lower and upper bound of each parameter. The conductivities are searched in
      \cmd{log10} space (\cmd{tx\_ost}), as the user guide recommends: a step from 1 to 10~m/d is as large as one from 10 to
      100.
\item \cmd{BeginResponseVars}: where to read each response: the file, the line (0 is the header line of
      \cmd{Diagnostics.csv}, so line 1 is the streamflow series), the column (5 = KGE, 4 = RMSE) and the separator. Check them
      against your own \cmd{Diagnostics.csv} (above) whenever you change the observations or the metrics.
\item \cmd{BeginTiedRespVars}: the objective \cmd{COST} $= -\mathrm{KGE} + \V{t6:w101}\,\mathrm{RMSE}_{101} +
      \V{t6:w102}\,\mathrm{RMSE}_{102} + \V{t6:w103}\,\mathrm{RMSE}_{103}$, with the RMSE in metres. Each weight is
      $1/(2\sigma)$, $\sigma$ the standard deviation of that well's observed heads in the calibration years: an error of
      half the well's own spread costs as much as a loss of 0.25 of KGE. For the initial guess, \cmd{COST} =
      \V{t6:costInit}.
\item \cmd{RandomSeed} makes the search repeatable; \cmd{MaxIterations} is the number of model runs;
      \cmd{UseInitialParamValues} starts from the \cmd{init} column (otherwise from random values).
\end{itemize}

\subsection{The scripts}
\showfile{t6_calibration/start/Ost-RAVEN.sh}
\cmd{\$RAVEN} is the variable of Tutorial~1: it must be exported in the terminal where OSTRICH runs. Removing
\cmd{Diagnostics.csv} before the run makes a failed run stop OSTRICH, instead of OSTRICH reading the results of the previous
run.
\showfile{t6_calibration/start/save_best.sh}
Both scripts must be executable (\cmd{chmod +x *.sh}; they are, in the tutorial files).

\section{Run DDS}
\begin{shell}
$ $OSTRICH
\end{shell}
OSTRICH prints little; it writes \cmd{OstModel0.txt} (one line per run: the objective and the parameters, in run order) and,
at the end, \cmd{OstOutput0.txt}. Follow the progress in another terminal with \cmd{tail -f OstModel0.txt}. The
\V{t6:maxiter} runs take about \V{t6:ddsmin}~minutes (about \V{t6:runsec}~seconds each). The end of \cmd{OstOutput0.txt}:
\term{t6_ostout.txt}
and its run record (each line a new best):
\term{t6_runrecord.txt}
\fig{t6_convergence.png}{DDS: the objective of each run and the best so far (left; grey: the best so far of the first
attempt), and the parameters of the best run so far divided by their true values (right).\label{fig:t6conv}}
\begin{tocheck}
\begin{itemize}
\item \cmd{OstErrors0.txt} reports \cmd{A total of 0 errors}, and \cmd{OstModel0.txt} has a header line and \V{t6:nevals} lines, one per run.
\item The objective fell from \V{t6:costInit} to \V{t6:costBest}; DDS had made 95\% of that improvement after
      \V{t6:itninetyfive} runs, and found its best at run \V{t6:itBest} (Figure~\ref{fig:t6conv}).
\item \cmd{best/} holds the \cmd{Creek.rvp}, \cmd{Diagnostics.csv}, \cmd{Hydrographs.csv} and \cmd{GWHeads.csv} of the best run.
\end{itemize}
\end{tocheck}
\begin{trouble}
\begin{itemize}
\item OSTRICH stops at once with an error about \cmd{Diagnostics.csv}: the first model run failed. Run
      \cmd{./Ost-RAVEN.sh} by hand and read \cmd{model/output/Raven\_errors.txt}; most often \cmd{\$RAVEN} is not exported.
\item Every run gives the same objective: the template's parameter names do not match \cmd{BeginParams}, so \cmd{Creek.rvp}
      never changes.
\item A response of \cmd{-1e+99} or \cmd{inf}: a series without observations in the evaluated period (check the weights).
\end{itemize}
\end{trouble}

\section{What DDS recovered}
\begin{center}\small
\begin{tabular}{lrrrrr}\toprule
Parameter & range & initial & true & calibrated & error\\\midrule
$K_h$ of SAND [m/d] & \V{t6:lowK_SAND}--\V{t6:highK_SAND} & \V{t6:initK_SAND} & \V{t6:trueK_SAND} & \V{t6:bestK_SAND} & \V{t6:errK_SAND}\%\\
$K_h$ of TILL [m/d] & \V{t6:lowK_TILL}--\V{t6:highK_TILL} & \V{t6:initK_TILL} & \V{t6:trueK_TILL} & \V{t6:bestK_TILL} & \V{t6:errK_TILL}\%\\\bottomrule
\end{tabular}\end{center}
The till conductivity is recovered within \V{t6:errK_TILL}\%. The sand conductivity ends \V{t6:errK_SAND}\% from its
true value, yet the objective of the calibrated set, \V{t6:costBest}, is no worse than that of the true parameters,
\V{t6:costcalibtrue}: with this noise the observations cannot tell a sand conductivity of \V{t6:bestK_SAND}~m/d from one of
\V{t6:trueK_SAND}~m/d (Section~\ref{sec:t6ident}: only OW\_SLOPE sees the sand, by centimetres, and its weight is small).
This is the identifiability limit of the observations, not of the search: more runs would not change it. An observation
that depends on the sand would: a pumping test (the transient drawdown when PW\_1 starts or stops pumping), or a well
where the sand carries the water of the slopes down the valley. Another well in the sand near PW\_1 would not help:
OW\_NEAR is one, and its seasonal heads do not see the sand.

\section{Validate on 2013--2014}
A calibrated model must also work on data it has not seen. Make a validation folder from \cmd{model/} with the best
\cmd{Creek.rvp}, run it for the five years 2010--2014, and let Raven score the calibration and validation years separately:
\begin{shell}
$ cp -r model validation
$ cp best/Creek.rvp validation/
$ cd validation
\end{shell}
\diff{t6_valid_rvi.diff}
\term{t6_run_valid.txt}
\term{t6_diag_valid.txt}
\begin{center}\small
\begin{tabular}{lrrrrrr}\toprule
 & \multicolumn{3}{c}{calibration years 2011--2012} & \multicolumn{3}{c}{validation years 2013--2014}\\
Parameters & initial & calibrated & true & initial & calibrated & true\\\midrule
KGE of streamflow & \V{t6:kgecalibinit} & \V{t6:kgecalib} & \V{t6:kgecalibtrue} & \V{t6:kgevalidinit} & \V{t6:kgevalid} & \V{t6:kgevalidtrue}\\
RMSE OW\_NEAR [m] & \V{t6:rmsecalib101init} & \V{t6:rmsecalib101} & \V{t6:rmsecalib101true} & \V{t6:rmsevalid101init} & \V{t6:rmsevalid101} & \V{t6:rmsevalid101true}\\
RMSE OW\_UPPER [m] & \V{t6:rmsecalib102init} & \V{t6:rmsecalib102} & \V{t6:rmsecalib102true} & \V{t6:rmsevalid102init} & \V{t6:rmsevalid102} & \V{t6:rmsevalid102true}\\
RMSE OW\_SLOPE [m] & \V{t6:rmsecalib103init} & \V{t6:rmsecalib103} & \V{t6:rmsecalib103true} & \V{t6:rmsevalid103init} & \V{t6:rmsevalid103} & \V{t6:rmsevalid103true}\\
\cmd{COST} & \V{t6:costcalibinit} & \V{t6:costcalib} & \V{t6:costcalibtrue} & \V{t6:costvalidinit} & \V{t6:costvalid} & \V{t6:costvalidtrue}\\\bottomrule
\end{tabular}\end{center}
(The columns \emph{initial} and \emph{true} are the same validation run with the initial and the true \cmd{Creek.rvp}.)
\begin{tocheck}
Compare the columns: the calibrated set is far better than the initial guess everywhere and as good as the true parameters
on the validation years too (KGE \V{t6:kgevalid} against \V{t6:kgevalidtrue}; head errors at the \V{t6:noiseh}-cm noise of the
observations at all three wells), although its sand conductivity is off: where the observations cannot see a parameter,
its error does no harm to what they measure --- and no good to predictions that depend on it.
\end{tocheck}
\fig{t6_validation.png}{Observed (grey) and simulated streamflow and heads with the initial (orange) and the calibrated
(blue) parameters; shaded: the warm-up year, dashed: the start of the validation years.}

\section{Going further}
\begin{itemize}
\item More parameters ($S_y$ of the till, \cmd{SEEPAGE\_LEAKANCE}, soil parameters): check first, as in Section~\ref{sec:t6ident}, that the
      observations see each of them, and on its own; DDS scales its search to the number of runs it is given.
\item Several processors: OSTRICH's parallel DDS (\cmd{ParallelDDS}, Demo~C3) runs each model in its own folder
      (\cmd{ModelSubdir}, \cmd{BeginExtraDirs}); every Raven run builds its own MODFLOW model, so parallel runs do not interfere.
\item A soil-zone model (Tutorial~3) calibrates the same way; its runs take longer.
\end{itemize}
The finished state, \cmd{\$TUT/t6\_calibration/finished}, holds \cmd{OstOutput0.txt}, \cmd{OstModel0.txt}, \cmd{best/} and
\cmd{validation/} of the calibration shown here; to keep it small, its \cmd{best/} holds only \cmd{Creek.rvp} and
\cmd{Diagnostics.csv} of the best run.
